% =====================================================================
% JMIR Research Protocols requires a structured abstract:
% Background / Objective / Methods / Results / Conclusions
% plus the International Registered Report Identifier (IRRID) line at the end.
% Length: 450 words max.
% =====================================================================

\begin{abstract}
\noindent \singlespacing

\noindent \textbf{Background:} Objective assessment of cognitive performance during ordinary teaching remains a methodological gap in educational research. Conventional instruments---summative tests, self-report, lesson-end ratings---operate at coarse temporal grains and carry well-documented subjective biases. Multimodal passive sensing offers continuous, ecologically valid monitoring through ambient, physiological, neurophysiological and behavioural signals, yet these technologies have matured in disconnected silos rather than as integrated classroom systems.

\noindent \textbf{Objective:} This protocol describes an observational, longitudinal study in two groups of Spanish compulsory secondary education students (ages 14--16) that integrates, in a four-month deployment, a multimodal passive sensing platform with a teacher-annotated cognitive-performance label and a Psychomotor Vigilance Task (PVT) external anchor. The primary objective is to characterise the platform's measurement properties before any inferential modelling.

\noindent \textbf{Methods:} Four passive streams are recorded in synchrony: (i)~a custom ambient platform (temperature, humidity, CO\textsubscript{2}, PM2.5, PM10, illuminance, sound pressure level at 15~Hz); (ii)~EmotiBit wristbands (photoplethysmography at 25~Hz, electrodermal activity at 15~Hz, skin temperature, and 9-axis inertial measurement at 100~Hz, six units rotating weekly); (iii)~a 32-channel semi-dry EEG system at 256~Hz on one participant per week; and (iv)~overhead 4K video processed exclusively for body-pose features, with raw footage deleted within 72~hours. The cognitive-performance label is a binary teacher annotation per 5-minute window. External validation uses the PVT administered monthly and an end-of-session 0--10 rating. The technical contribution is a dual-plane architecture using MQTT QoS~2 and NTP timestamping, targeting sub-100~ms inter-modality alignment. The study is observational; no intervention is delivered.

\noindent \textbf{Results:} The protocol has been submitted for ethics evaluation (CEIS-UCLM, reference pending). The acquisition platform is implemented and bench-tested. Upon favourable evaluation, the consent-and-assent campaign will open; data collection is expected to begin in the subsequent school term and to run for four consecutive months. Results are anticipated within 12 months of deployment completion.

\noindent \textbf{Conclusions:} The protocol establishes a measurement infrastructure that will support subsequent predictive analyses within the doctoral compendium. The joint operationalisation of four passive multimodal streams with a teacher-anchored cognitive-performance label and a PVT external benchmark, under European data-protection guarantees for minors, provides a replicable framework for future multimodal learning analytics deployments in compulsory secondary education.

\end{abstract}
